\chapter{Introduction}
\label{ch:introduction}

\section{Purpose and Scope}

The SIRC-1 CPU is sold on its own, so it is documented in three manuals that stand apart from any system built around
it. SIRCADA and SIRCAS are made for the SIRCUS system and each has a manual of its own. Together the five manuals have
the following scopes:

\begin{table}[H]
	\centering
	\begin{tabular}{lp{8.2cm}}
		\toprule
		\textbf{Manual} & \textbf{Scope}                                                                                                                \\
		\midrule
		\archmanual{}   & ISA, instruction encodings, canonical assembly notation, exceptions, programmer-visible state                                 \\
		\asmmanual{}    & Complete assembly syntax, directives, symbols, meta-instructions, and toolchain usage                                         \\
		\hwmanual{}     & Pinout, bus interface, timing, electrical characteristics                                                                     \\
		\videomanual{}  & The display adapter: package and pinout, memory map, registers, data formats, video modes, timing                             \\
		\audiomanual{}  & The sound processor: package and pinout, memory map, registers, data formats, voices, filter, sequencer, timing (this manual) \\
		\bottomrule
	\end{tabular}
	\caption{The SIRC-1 and SIRCUS Manual Set}
	\label{tab:manual-set}
\end{table}

This manual covers SIRCAS as a component: its package and pins, what the CPU can read and write, how the device turns
the contents of its audio memory into sound, and the timing of its signals. It does not describe the CPU, its
instructions or its bus protocol. The device is a bus peripheral like any other, so the \hwmanual{} is authoritative
for how the CPU makes a bus cycle, and this manual is authoritative for how the device answers it and for what it does
once it has been reached.

\section{Device Overview}

SIRCAS plays six voices at once. Each voice reads eight-bit samples from the device's own audio memory at a pitch set
by a note number, shapes them with an envelope and a volume, and places them in a stereo mix, either directly or
through a resonant filter shared by all the voices. Six small sequencers, one for each voice, can play a tune from
instructions in audio memory with no help from the CPU. The mixed result leaves the device as a serial stream of signed
16-bit stereo samples for an external digital-to-analogue converter.

\begin{table}[H]
	\centering
	\begin{tabular}{lp{9cm}}
		\toprule
		\textbf{Property} & \textbf{Value}                                                                                                      \\
		\midrule
		Voices            & Six, each playing signed 8-bit samples by nearest-sample lookup (no interpolation)                                  \\
		Sample rate       & 31,960.226 Hz: one output sample every 336 \texttt{CLKI} cycles                                                     \\
		Audio memory      & 32,768 words (64 KB) of external static memory, dedicated to the device (Chapter~\ref{ch:package})                  \\
		Pitch             & 16.16 fixed-point step; A 12-entry pitch table for each instrument; notes 0--119 (Chapter~\ref{ch:playback})        \\
		Envelope          & 12-bit linear attack, decay, sustain and release, one step every 64 samples (Chapter~\ref{ch:envelope})             \\
		Mixer             & Per-voice volume (8 bits) and pan (8 bits), master volume, 20-bit accumulators (Chapter~\ref{ch:mixer})             \\
		Modulation        & Any voice can modulate pitch, volume, pan, filter cutoff, filter pan or master volume (Chapter~\ref{ch:modulation}) \\
		Filter            & One shared two-pole state-variable filter: low-pass, band-pass or high-pass (Chapter~\ref{ch:filter})               \\
		Sequencers        & Six, one per voice, with a shared tick; 16-bit instructions (Chapter~\ref{ch:sequencer})                            \\
		Output            & Stereo, signed 16 bits, serial (I\textsuperscript{2}S framing) to an external converter                             \\
		Interrupts        & One: sequence completion (\activelow{IRQ3})                                                                         \\
		Clock             & \texttt{CLKI} at 10,738,636 Hz (half the 21,477,272 Hz system oscillator)                                           \\
		\bottomrule
	\end{tabular}
	\caption{Sound Device Summary}
	\label{tab:device-summary}
\end{table}

\section{The Voice Signal Path}
\label{sec:signal-path}

Every voice follows the same path, shown in Figure~\ref{fig:signal-path}. The voice reads a sample byte from audio
memory at the position its position accumulator selects. The envelope and the voice's volume scale the sample. What
happens next depends on the voice's routing bits in the \texttt{CONTROL} register (Section~\ref{sec:reg-control}):

\begin{itemize}
	\item A \textbf{dry voice} is divided between the left and right accumulators by its pan value.
	\item A \textbf{filtered voice} is added to the input of the shared filter instead. The filter's output has a volume and a
	      pan of its own, and is added to the same accumulators.
	\item A \textbf{modulation voice} makes no sound. Its scaled sample steers a parameter of another voice or of the device
	      (Chapter~\ref{ch:modulation}).
\end{itemize}

The master volume scales the two accumulators, and the result is limited to 16 bits and sent to the converter.

\begin{figure}[H]
	\centering
	\begin{tikzpicture}[
			scale=0.9, transform shape,
			block/.style={draw, rounded corners=2pt, fill=sirclightgray, align=center, font=\footnotesize, minimum height=0.9cm, minimum width=2.5cm},
			store/.style={draw, fill=white, align=center, font=\footnotesize, minimum height=0.9cm, minimum width=2.5cm},
			flow/.style={-{Stealth[length=2mm]}, thick},
			ctl/.style={-{Stealth[length=2mm]}, thick, dashed},
			note/.style={font=\scriptsize, text=sircgray}
		]
		\node[store] (ram) at (0,0) {Audio memory\\(samples)};
		\node[block] (pcm) at (3.4,0) {Sample\\playback};
		\node[block] (env) at (6.6,0) {Envelope\\and volume};
		\node[block] (pan) at (9.8,1.2) {Pan\\(dry voice)};
		\node[block] (flt) at (9.8,-1.2) {Shared filter\\(filtered voice)};
		\node[block] (acc) at (13.0,0) {Left and right\\accumulators};
		\node[block] (mst) at (13.0,-2.2) {Master\\volume};
		\node[block] (out) at (13.0,-4.2) {Serial\\output};
		\node[block] (mod) at (6.6,2.6) {Modulation\\(voice in modulation mode)};
		\draw[flow] (ram) -- (pcm);
		\draw[flow] (pcm) -- (env);
		\draw[flow] (env.east) -- ++(0.5,0) |- (pan.west);
		\draw[flow] (env.east) -- ++(0.5,0) |- (flt.west);
		\draw[flow] (pan.east) -| ([xshift=-0.5cm]acc.north);
		\draw[flow] (flt.east) -| ([xshift=0.5cm]acc.south);
		\draw[flow] (acc) -- (mst);
		\draw[flow] (mst) -- (out);
		\draw[ctl] (mod.east) -| node[above, near start, note]{pitch, volume, pan, filter, master} (acc.north east);
	\end{tikzpicture}
	\caption{The Voice Signal Path}
	\label{fig:signal-path}
\end{figure}

\section{Block Diagrams}
\label{sec:block-diagrams}

Figure~\ref{fig:block-system} shows the device as the CPU sees it: a register block and a window onto audio memory
behind one bus interface, a serial output, and one interrupt line. The audio memory is not on the CPU's bus. It is a
separate memory behind the device, and the device alone controls it; the CPU reaches it only through the device. The
device and its memory are described in Chapters~\ref{ch:package} and~\ref{ch:memory-map}.

\begin{figure}[H]
	\centering
	\begin{tikzpicture}[
			scale=0.88, transform shape,
			block/.style={draw, rounded corners=2pt, fill=sirclightgray, align=center, font=\small},
			mem/.style={draw, fill=white, align=center, font=\footnotesize, minimum height=0.9cm, minimum width=3.4cm},
			flow/.style={-{Stealth[length=2mm]}, thick},
			note/.style={font=\scriptsize, text=sircgray}
		]
		\node[block, minimum width=2.2cm, minimum height=3.0cm] (cpu) at (0,0) {SIRC-1\\CPU};
		\node[block, minimum width=3.4cm, minimum height=3.0cm] (bus) at (4.4,0) {Bus interface\\and address\\decode};
		\node[mem] (m1) at (9.0,0.9) {Registers\\\texttt{0x0000}--\texttt{0x006F}};
		\node[mem] (m2) at (9.0,-0.9) {Audio memory window\\\texttt{0x8000}--\texttt{0xFFFF}};
		\node[block, minimum width=3.0cm, minimum height=3.0cm] (snd) at (13.2,0) {Sound\\processor};
		\node[mem] (sram) at (13.2,-3.4) {Audio memory\\32,768 words};
		\node[block, minimum width=3.0cm] (dac) at (13.2,3.4) {Serial output\\to converter};
		\draw[flow, <->] (cpu) -- node[above, note]{address, data} (bus);
		\draw[flow, <->] (bus.east) -- ++(0.6,0) |- (m1.west);
		\draw[flow, <->] (bus.east) -- ++(0.6,0) |- (m2.west);
		\draw[flow, <->] (m1.east) -- (m1.east -| snd.west);
		\draw[flow, <->] (m2.east) -- (m2.east -| snd.west);
		\draw[flow, <->] (snd) -- (sram);
		\draw[flow] (snd) -- (dac);
		\draw[flow] (snd.north west) -- ++(0,0.9) -| node[above, near start, note]{\activelow{IRQ3}} (cpu.north);
	\end{tikzpicture}
	\caption{The Sound Device in the System}
	\label{fig:block-system}
\end{figure}

Figure~\ref{fig:block-spu} shows what is inside. One slot scheduler decides, in a fixed pattern, whether the audio
memory is used by the sound processor or by the CPU's window (Chapter~\ref{ch:timing}). The sound processor uses its
slots for sample reads first, and for the sequencers' instruction, instrument and pitch-table reads with what is left.

\begin{figure}[H]
	\centering
	\begin{tikzpicture}[
			scale=0.86, transform shape,
			block/.style={draw, rounded corners=2pt, fill=sirclightgray, align=center, font=\footnotesize, minimum height=0.95cm, minimum width=3.1cm},
			store/.style={draw, fill=white, align=center, font=\footnotesize, minimum height=0.95cm, minimum width=3.1cm},
			flow/.style={-{Stealth[length=2mm]}, thick},
			ctl/.style={-{Stealth[length=2mm]}, thick, dashed},
			note/.style={font=\scriptsize, text=sircgray}
		]
		\node[block] (regs) at (0,6.6) {Register file\\(read by every block)};
		\node[block] (tick) at (4.4,6.6) {Tick generator\\\texttt{SEQ\_PERIOD}};
		\node[block] (seq) at (8.8,6.6) {Six sequencers\\one shared decoder};
		\node[store] (sram) at (0,3.2) {Audio memory};
		\node[block, minimum height=3.4cm] (arb) at (4.4,3.2) {Slot scheduler\\\scriptsize sound processor 1 slot,\\\scriptsize CPU window 3 slots};
		\node[block] (pcm) at (8.8,4.4) {Six sample\\players};
		\node[block] (env) at (8.8,2.4) {Six envelope\\generators};
		\node[block] (mod) at (8.8,0.4) {Modulation\\routing};
		\node[block] (mix) at (13.2,4.4) {Mixer\\gain, pan};
		\node[block] (flt) at (13.2,2.4) {Shared filter};
		\node[block] (mst) at (13.2,0.4) {Master volume\\and limiter};
		\node[block] (ser) at (13.2,-1.6) {Serial output};
		\node[block] (cpu) at (0,0.4) {CPU window\\and registers};
		\draw[ctl] (tick) -- (seq);
		\draw[flow, <->] (sram) -- (arb);
		\draw[flow, <->] (cpu) |- (arb.south west);
		\draw[flow] (arb.east |- pcm) -- (pcm.west);
		\draw[flow] (seq.south) -- ++(0,-0.4) -| (arb.north east);
		\draw[flow] (pcm.south) -- (env.north);
		\draw[flow] (env.east) -- ++(0.5,0) |- (mix.west);
		\draw[flow] (env.east) -- ++(0.5,0) |- (flt.west);
		\draw[flow] (mix.south) -- (flt.north);
		\draw[flow] (flt.south) -- (mst.north);
		\draw[flow] (mst.south) -- (ser.north);
		\draw[ctl] (mod.east) -- (mod.east -| mst.west);
	\end{tikzpicture}
	\caption{Blocks of the Sound Processor}
	\label{fig:block-spu}
\end{figure}

\section{Conventions}

\begin{itemize}
	\item Register and window addresses are in 16-bit words, as on the CPU. A word address is an offset from the base of the
	      device's segment unless stated otherwise. Sample data and some fields of the instrument table are byte offsets into
	      audio memory; each is marked where it occurs (Section~\ref{sec:byte-offsets}).
	\item Bit 0 is the least significant bit. Bit diagrams draw the most significant bit on the left.
	\item Hexadecimal values carry a \texttt{0x} prefix and upper case digits.
	\item Ranges of values and addresses use an en dash: \texttt{0x0000}--\texttt{0x00FF}.
	\item \textbf{Reserved} fields and registers are set aside by the architecture. Software must not rely on what they
	      read, and should write zero to a reserved field it does not otherwise use.
	\item A value that the design has not fixed is \textbf{not yet specified}. It is not undefined, and no value is implied.
	\item Fixed-point numbers are written $m.n$: $m$ integer bits and $n$ fraction bits. The position and step of a voice are
	      16.16 values, so their integer part is the high 16 bits.
	\item The \textbf{device} is the whole package, SIRCAS. It contains the \textbf{sound processor}, which is the six voices,
	      the mixer and the filter, together with the registers, the sequencers and the memory slot scheduler.
	\item A \textbf{tick} is one step of the sequencers' clock, whose period is set by \texttt{SEQ\_PERIOD}
	      (Chapter~\ref{ch:sequencer}). An \textbf{envelope update} is one step of a voice's envelope, made once in every 64
	      output samples (Chapter~\ref{ch:envelope}); portamento steps on it too.
	\item \textbf{Architecturally undefined} means that the manual does not say what the device does: the result may differ
	      between devices and between runs, and software must not depend on it. It is not the same as \textbf{not yet
		      specified}, which marks a decision the design has still to make, nor as the contents of audio memory at reset, which
	      are \textbf{undefined} until the program writes them.
	\item In formulas, $\ll$ and $\gg$ are shifts left and right (a right shift of a negative number rounds towards minus
	      infinity), $\lfloor x \rfloor$ rounds down, $\lceil x \rceil$ rounds up, and $a \bmod b$ is the remainder of $a$
	      divided by $b$.
\end{itemize}
